\section{Giới hạn và hướng mở rộng}
\label{sec:gioi-han}

\subsection{Giới hạn của các kết quả đã có}

\subsubsection{Phần A}

\begin{itemize}[leftmargin=2em]
\item \textbf{PPL của A.2 không so sánh được ra ngoài.} Giá trị $10{,}1948$ gắn với
một corpus cụ thể, một từ vựng ký tự 133 phần tử, một split cố định và một quy ước
chấm token. Nó không so được với PPL của mô hình khác dùng tokenizer khác. So sánh
perplexity chỉ có nghĩa khi cùng tokenizer, cùng dữ liệu và cùng quy ước chấm
(Mục~\ref{sec:b6-agg}).
\item \textbf{Độ lệch chuẩn $0{,}0628$ đo biến thiên hẹp.} Nó đo biến thiên giữa
các lần huấn luyện \emph{trên cùng một split}. Nó không bao gồm biến thiên do chọn
split khác, và không phải khoảng tin cậy.
\item \textbf{Chưa đối chiếu môi trường của A.2.} Thiếu log PyTorch/CUDA và
checkpoint, nên hạng mục tái lập của A.2 chưa đóng được
(Mục~\ref{sec:tai-lap}).
\item \textbf{Mô hình A.2 rất nhỏ} (2 khối, $\dmodel=64$). Nó là bằng chứng rằng
cài đặt đúng và huấn luyện được, không phải một mô hình ngôn ngữ hữu dụng.
\end{itemize}

\subsubsection{Phần B}

\begin{itemize}[leftmargin=2em]
\item \textbf{B.1 dựa trên một mô hình giả định.} Giả thiết độc lập, kỳ vọng $0$,
phương sai $1$ mô tả trạng thái tại khởi tạo, không phải phân phối của $q,k$ trong
mô hình đã huấn luyện. Kết quả giải thích lý do chọn thang điểm, không bảo đảm
không bão hòa.
\item \textbf{B.3 chỉ đúng với số chiều chẵn}, và Mục~\ref{sec:b3-odd} chứng minh
mệnh đề \emph{thất bại} ở chiều lẻ --- mà cài đặt của nhóm lại chấp nhận chiều lẻ.
Thí nghiệm hiện tại dùng chiều chẵn nên không chạm phải, nhưng đây là một lỗ hổng
thật giữa chứng minh và mã, không phải chi tiết hình thức.
\item \textbf{B.4 nói về chi phí số học, không nói về độ trễ.} Kết luận hợp nhất
không bảo đảm khớp từng bit trong số học dấu phẩy động, không tính chi phí tạo
$W'$, và không mở rộng thành kết luận về tốc độ khi đổi kernel hay tái lượng tử
hóa.
\item \textbf{B.5 chứng minh tính duy nhất của $\pi^*$, không phải của $\theta$.}
Việc mất mát DPO có nghiệm $\theta$ duy nhất hay không là câu hỏi khác, và không
được suy ra từ Mục~\ref{sec:b5}. Với miền đáp án không bị chặn reward, cần kiểm
riêng việc chuẩn hóa.
\end{itemize}

\subsubsection{Phần C}

\begin{itemize}[leftmargin=2em]
\item \textbf{C.1 và C.4 dùng đầu vào minh họa}, không phải bộ số giảng viên cấp.
Nhóm đã ghi rõ nguồn và cách tái lập, nhưng nếu nhận được đầu vào chính thức thì
phải thay số.
\item \textbf{C.2 chỉ đếm MAC của các tích ma trận} trên \emph{một} khối,
\emph{một} chuỗi, \emph{forward}. Không phải FLOP huấn luyện toàn mô hình. Các thao
tác không phải matmul (LayerNorm, softmax, ReLU, cộng bias) được nêu bằng
$\Theta(\cdot)$ chứ không cộng vào tổng.
\item \textbf{C.3 là ước lượng bộ nhớ, không phải phép đo.} Chưa cộng tái tính
activation, buffer giải lượng tử, workspace GEMM, CUDA context, tensor tạm của
optimizer và phân mảnh allocator. \textbf{Không} được kết luận từ bảng rằng một GPU
cụ thể chạy vừa SFT hay DPO.
\end{itemize}

\subsubsection{Phần D}

Giới hạn lớn nhất của báo cáo nằm ở đây: Phần~D chưa có kết quả thực nghiệm. Chi
tiết ở Mục~\ref{sec:phan-d}; tóm lại là D.2 còn 24/2.000 QA, D.3 chưa huấn luyện,
D.4 và D.5 chưa bắt đầu, và cầu nối A$\leftrightarrow$D (G5) chưa đóng.

\subsection{Rủi ro thiết kế đã nhận diện}

\subsubsection{Kiến thức thủ tục thay đổi nhanh}

Đây là rủi ro cấu trúc, không phải lỗi triển khai. Hai mâu thuẫn tìm được ở
Mục~\ref{sec:d1} đã minh họa: trong vòng chưa đầy một năm, thời hạn cấp trích lục
được quy định khác nhau ở hai quyết định, và nhóm không tìm được văn bản triển khai
tại Hà Nội để kết luận văn bản nào đang áp dụng.

Một mô hình được tinh chỉnh \emph{đóng băng} kiến thức tại thời điểm huấn luyện.
Khi văn bản thay đổi, mô hình trả lời sai mà không có cách nào tự biết. Ba cách
giảm nhẹ nhóm áp dụng trong phạm vi hiện tại:

\begin{enumerate}[leftmargin=2em]
\item Buộc câu trả lời dẫn số hiệu văn bản và phiên bản, để người dùng tự thấy đáp
án gắn với văn bản nào.
\item Soạn QA gắn với \emph{từng văn bản và phạm vi của nó}, thay vì phát biểu
chung ``hiện hành''.
\item Dạy mô hình từ chối khi chưa đủ căn cứ, thay vì nội suy.
\end{enumerate}

\subsubsection{Rủi ro về nhãn ưu tiên}

Nhãn ưu tiên là tín hiệu học duy nhất của DPO, nên chất lượng nhãn quyết định chất
lượng mô hình. Ba rủi ro:

\begin{itemize}[leftmargin=2em]
\item \textbf{Hai người chấm có thể lệch hệ thống} (ví dụ một người khắt khe hơn về
trích dẫn). Vì vậy phải \emph{đo} mức đồng thuận, không chỉ phân xử.
\item \textbf{Phân xử trước khi khóa nhãn sẽ làm hỏng phép đo đồng thuận.} Quy
trình đặt bước khóa trước bước phân xử chính vì lý do này.
\item \textbf{Cặp phải cùng prompt.} Theo \eqref{eq:b5-invert}, hạng
$\beta\log Z(x)$ chỉ triệt tiêu khi hai đáp án thuộc cùng một câu hỏi. Ghép đáp án
của hai câu hỏi khác nhau làm sai mô hình xác suất, không chỉ làm dữ liệu kém.
\end{itemize}

\subsubsection{Rủi ro đánh giá}

Nếu người đánh giá biết đáp án nào thuộc mô hình nào, win rate mất ý nghĩa. Nếu
người chấm nhãn D.4 cũng là người đánh giá D.5 và đã xem tập test, phép đo bị nhiễm.
Đây là lý do nhóm tách quyền truy cập \texttt{USER\_TEST}
(Mục~\ref{sec:dao-duc}).

\subsection{Hướng mở rộng}

\subsubsection{Truy xuất tăng cường (RAG)}

Hướng tự nhiên nhất cho vấn đề kiến thức thay đổi nhanh là truy xuất tăng cường:
thay vì ghi nội dung thủ tục vào trọng số, cho mô hình truy xuất văn bản tại thời
điểm trả lời và sinh đáp án dựa trên đoạn đã truy xuất. Cách này có hai lợi thế
đúng với bài toán: cập nhật văn bản không cần huấn luyện lại, và trích dẫn trở thành
hệ quả tự nhiên của kiến trúc chứ không phải hành vi phải dạy.

Đây là Đề tài 11 trong danh mục của đề. Nhóm \textbf{chỉ thảo luận} đây là hướng mở
rộng và \textbf{không} trình bày RAG như hạng mục đã triển khai; mở rộng sang RAG
không nằm trong phạm vi Đề tài 9.

\subsubsection{Các mở rộng khác trong phạm vi Đề tài 9}

\begin{itemize}[leftmargin=2em]
\item \textbf{Vá B.3 trong mã:} buộc $\dmodel$ chẵn, hoặc bỏ tọa độ $\sin$ lẻ cặp,
để tính chất dịch vị trí đúng với mọi cấu hình mà mã cho phép.
\item \textbf{Mở rộng adapter ngoài Q/V:} C.3 tính cho Q/V theo đề. Gắn thêm
adapter vào K, O hoặc các ma trận FFN làm tăng số tham số học theo
\eqref{eq:b4-count} và có thể đổi kết quả; đây là câu hỏi thực nghiệm chưa trả lời.
\item \textbf{Đo peak VRAM thật} để thay ước lượng C.3 bằng số đo, ghi cả peak
allocated và peak reserved.
\item \textbf{Mở rộng tập test khảo sát} sang nhiều nhóm người dùng, vì cách diễn
đạt khác nhau giữa các nhóm tuổi và mức quen thuộc với thủ tục hành chính.
\end{itemize}

\subsection{Nhóm không hứa điều gì}

Để tránh ngộ nhận, nhóm nêu rõ ba điều \emph{không} được suy ra từ báo cáo này:

\begin{enumerate}[leftmargin=2em]
\item \textbf{Không hứa DPO tốt hơn SFT.} Giả thuyết ở Mục~\ref{sec:d5} là giả
thuyết. Nhóm sẽ báo cả kết quả âm nếu đo được.
\item \textbf{Không coi bằng chứng kỹ thuật là bằng chứng quyền.} Việc audit chạy
được không có nghĩa nhóm có quyền huấn luyện hay phân phối dữ liệu đó.
\item \textbf{Không coi kế hoạch là kết quả.} Mọi hạng mục đã lập kế hoạch nhưng
chưa chạy đều mang nhãn \chualam.
\end{enumerate}
